\section{Sampling an original construction}
\label{sec:peps_original_compression_sampling}

We now combine the local approximation, pair-effect replacement and source
sampling estimates for the original gate maps. These are the approximation
steps in \cite[proof of Theorem~5.2]{OpenAI2026PolynomialPEPS}.
% Source: 04-compression.tex, lines 137--175, 199--229 and 590--600.
The sampled physical operator need not be positive. Its identification with a
contraction of local tensors is a separate assertion.

\begin{definition}[Distributed construction]
    \label{def:peps_distributed_construction}
    \lean{TNLean.PEPS.PairEffect.DistributedConstruction}
    \leanok
    \uses{def:peps_original_circuit}
    A distributed construction has a finite nonempty set $P$ of parties.
    For each $p\in P$, it specifies a finite-dimensional initial private
    memory $H_p$ and a unit vector $x_p\in H_p$. It also specifies a physical
    dimension $d_p$, an ordering of the parties without repetition, a finite
    ordered collection $R$ of finite-dimensional private output registers,
    and an actual contraction circuit
    \begin{align}
        W:\bigotimes_{p\in P}H_p
        \longrightarrow
        \left(\bigotimes_{p\in P}\mathbb C^{d_p}\right)\otimes\mathcal H_R.
    \end{align}
    The circuit retains the original gate expansions and participant labels.
    There is no numerical bound on private dimensions in this definition.
    Resource estimates and approximation conclusions are separate statements.

    The nonempty-party convention permits scalar factors to be absorbed in
    an existing local tensor, as in the proof of Theorem~5.2. It does not
    exclude gates with no participating parties. The zero-party scalar
    boundary case is recorded in \cite{gap:polypeps_nonempty_parties}.
\end{definition}

\begin{definition}[Original input and physical density]
    \label{def:peps_distributed_construction_density}
    \lean{TNLean.PEPS.PairEffect.DistributedConstruction.input,
        TNLean.PEPS.PairEffect.DistributedConstruction.physicalDensity}
    \leanok
    \uses{def:peps_distributed_construction}
    The input of a distributed construction is the normalized product
    $\psi=\bigotimes_{p\in P}x_p$. Its physical density is
    \begin{align}
        \rho=\operatorname{Tr}_{\mathcal H_R}
        \bigl(|W\psi\rangle\langle W\psi|\bigr),
    \end{align}
    expressed in the physical coordinates indexed by the original parties.
    The private trace uses a finite orthonormal basis of the actual retained
    registers. No normalization by $\|W\psi\|$ is performed.
\end{definition}

\begin{theorem}[Subnormalized physical density]
    \label{thm:peps_distributed_construction_density}
    \lean{TNLean.PEPS.PairEffect.DistributedConstruction.physicalDensity_posSemidef,
        TNLean.PEPS.PairEffect.DistributedConstruction.trace_physicalDensity_le_one}
    \leanok
    \uses{def:peps_distributed_construction_density}
    The actual physical density of a distributed construction is positive
    semidefinite and satisfies $\operatorname{Tr}\rho\leq1$.
\end{theorem}
\begin{proof}
    \leanok
    The outer product of the original output vector is positive semidefinite,
    and partial trace preserves positivity. Partial trace also preserves the
    full trace, so $\operatorname{Tr}\rho=\|W\psi\|^2$. Each gate is a
    contraction and $\|\psi\|=1$, giving the asserted bound.
\end{proof}

\begin{theorem}[A fixed inverse power meets the actual gate budget]
    \label{thm:peps_fixed_power_gate_budget}
    \lean{TNLean.PEPS.PairEffect.inversePower_le_sourceGateBudget_half}
    \leanok
    Let $L\ge2$, $C\ge1$, $M\le CL^n$, where $M,n$ are nonnegative
    integers, and let $p$ be real. Set
    $q=n+\lceil p\rceil_++\lceil16C\rceil_+$. Then
    \begin{align}
        L^{-q}\le \frac{L^{-p}}{16\max\{1,M\}}.
    \end{align}
    The exponent is fixed by the original count bound and the target accuracy,
    and works simultaneously at every size $L\ge2$.
\end{theorem}
\begin{proof}
    \leanok
    \uses{thm:peps_approximate_gate_real_accuracy}
    The bound $\max\{1,M\}\le CL^n$ follows from $C,L\ge1$.
    Apply Theorem~\ref{thm:peps_approximate_gate_real_accuracy}
    with this enlarged count,
    and divide by its positive denominator.
\end{proof}

\begin{theorem}[Using the original local approximation budget]
    \label{thm:peps_original_approximation_budget}
    \lean{TNLean.PEPS.PairEffect.EffectCircuit.IsGateApproximation.mono}
    \lean{TNLean.PEPS.PairEffect.EffectCircuit.isGateApproximation_budget_of_power_bound}
    \leanok
    Increasing the permitted local approximation error preserves the original
    operators, approximate monomial lists and all contraction and allowedness
    conditions. Consequently, under the preceding count assumptions, gate
    expansions available at accuracy $L^{-q}$ satisfy the local budget
    $L^{-p}/(16\max\{1,M\})$ required to reserve one quarter of the final
    physical error for this approximation step.
\end{theorem}
\begin{proof}
    \leanok
    \uses{thm:peps_fixed_power_gate_budget}
    Only the inequality bounding each local operator difference changes.
    Apply Theorem~\ref{thm:peps_fixed_power_gate_budget} and transitivity
    at every original occurrence. All the other local hypotheses are unchanged.
\end{proof}

\begin{theorem}[A positive polynomial coefficient bound]
    \label{thm:peps_original_coefficient_majorant}
    \lean{TNLean.PEPS.PairEffect.OriginalCircuit.exists_one_le_isExpansionBounded_of_termwise_power_bounds}
    \leanok
    Let $L\geq1$. Suppose every nonprivate gate of the original circuit
    has at most $K\leq C_KL^u$ monomials, each with at most $r$ pair
    effects and coefficient of absolute value at most $C_CL^v$, where
    $C_K,C_C\geq0$ and $u,v$ are nonnegative integers. Then there is
    a bound $B\geq1$ for the absolute coefficient sum of every gate,
    with the same effect bound $r$, such that
    \begin{align}
        B\leq (1+C_KC_C)L^{u+v}.
    \end{align}
    Thus the lower bound $B\geq1$ required in the source-sampling
    estimate follows from the original polynomial assumptions.
\end{theorem}
\begin{proof}
    \leanok
    \uses{thm:peps_original_coefficient_sum}
    Choose $B=(1+C_KC_C)L^{u+v}$. This is at least one since $L\geq1$.
    Each coefficient sum is at most
    $K C_CL^v\leq C_KC_CL^{u+v}\leq B$.
\end{proof}

\begin{theorem}[Finite expansion bounds and replacement coefficients]
    \label{thm:peps_original_finite_expansion_bound}
    \lean{TNLean.PEPS.PairEffect.OriginalCircuit.isMonomialBounded_mono}
    \lean{TNLean.PEPS.PairEffect.OriginalCircuit.exists_isMonomialBounded}
    \lean{TNLean.PEPS.PairEffect.OriginalCircuit.sum_norm_branchCoefficient_physicalFirstReplacement_le}
    \leanok
    Every finite original circuit has an integer $K$ bounding the lengths
    of all its gate expansions. Any larger integer is also such a bound.
    Suppose its monomials contain at most $r$ pair effects and the absolute
    coefficient sum at every gate is at most $B$. For any positive
    replacement accuracy, every gate of the actual source replacement,
    with the physical output registers placed first, has absolute branch
    coefficient sum at most $B$.
\end{theorem}
\begin{proof}
    \uses{thm:peps_physical_first_occurrences,thm:peps_source_only_resources}
    \leanok
    Take a common bound on the finitely many original gate lists. The
    coefficient estimate of Theorem~\ref{thm:peps_physical_first_occurrences}
    then gives the bound. The final exchange of physical and auxiliary registers
    introduces no gate and changes no coefficient.
\end{proof}

\begin{theorem}[Physical approximation after replacement and sampling]
    \label{thm:peps_sampled_original_density}
    \lean{TNLean.PEPS.PairEffect.OriginalCircuit.exists_physicalFirst_sourceSample_original_error_le}
    \leanok
    Let $W$ be a circuit on a finite party set $P$, with initial vector
    $\psi=\bigotimes_{p\in P}\psi_p$, where each $\psi_p$ is normalized
    in the whole initial private memory of party $p$. The final physical
    register at $p$ has dimension $d_p$, and every discarded private
    register is finite-dimensional. Suppose that the original nonprivate
    gate count is at most $M$, each party participates in at most $\ell$
    nonprivate occurrences, and the original expansion bounds hold with
    effect budget $r$ and coefficient-sum bound $B\ge1$. Let $D\ge1$
    satisfy $d_p\le D$ for every $p$.

    Fix finite orthonormal coordinates on the whole initial memory.
    Fix $0<\varepsilon\le1$. Set
    $\delta=\varepsilon/(8\max\{1,M\})$, and form the physical-first
    replacement $R$ using the corresponding prepared stack length.
    Let $\mathcal E_R$ be its actual source occurrences and
    $N_R=|\mathcal E_R|$. Choose
    \begin{align}
        Q & =B^{4\ell}D^4, &
        k & =\left\lceil
        \left(\frac{8\max\{1,N_R\}Q}{\varepsilon}\right)^2
        \right\rceil+1.
        \label{eq:peps_original_sample_count}
    \end{align}
    There exist Schmidt data for every actual source and branch, and
    one Gaussian sample array $\omega$, such that the sampled physical
    operator $\widehat\rho_R(\omega)$ satisfies
    \begin{align}
        \|\widehat\rho_R(\omega)-\rho_W\|_1
        \le\frac{3\varepsilon}{4}.
        \label{eq:peps_original_sample_error}
    \end{align}
    Here $\rho_W$ is the physical partial trace of $|W\psi\rangle
        \langle W\psi|$, in the original party-labelled coordinates.
    No numerical bound on private-memory dimensions is assumed.
    These are the density-approximation steps in
    \cite[proof of Theorem~5.2]{OpenAI2026PolynomialPEPS}.
    % Source: 04-compression.tex, lines 137--151, 199--229 and 342--381.
\end{theorem}

\begin{proof}\leanok
    \uses{thm:peps_physical_first_exact_sources,thm:peps_actual_source_sampling,
        thm:peps_auxiliary_register_finiteness,thm:peps_original_finite_expansion_bound}
    Prepare each $\psi_p$ by a local contraction from the scalar space.
    Their chronological composition prepares $\psi$ and introduces no
    pair source. The replacement is a contraction, preserves the original
    participating occurrences, and retains coefficient-sum bound $B$.
    A finite bound on its monomial counts follows from the finite original
    gate lists; it is not an additional hypothesis. Every auxiliary
    register is finite-dimensional by its Euclidean construction, so the
    individual private-register assumption suffices for all discarded
    registers of $R$.

    Apply Theorem~\ref{thm:peps_actual_source_sampling} to this actual
    circuit, preparation, and discarded memory. It gives original-source
    Schmidt data and one sample $\omega$ with
    $\|\widehat\rho_R(\omega)-\rho_R\|_1\le\varepsilon/4$.
    Theorem~\ref{thm:peps_physical_first_exact_sources} gives
    $\|\rho_R-\rho_W\|_1\le\varepsilon/2$. The trace-norm triangle
    inequality yields~(\ref{eq:peps_original_sample_error}).
\end{proof}


\begin{theorem}[Sampling from approximate original gates]
    \label{thm:peps_approximate_source_sample}
    \lean{TNLean.PEPS.PairEffect.EffectCircuit.exists_sourceSample_approximate_error_le}
    \leanok
    Let a finite distributed construction begin in a product of normalized
    party vectors. Suppose its original gates are contractions, its number of
    nonprivate gates is at most $M$, and each party participates in at most $b$
    such gates. All private spaces are finite-dimensional, without numerical
    dimension bounds. Let the physical dimension at each party be at most
    $D\ge1$.

    Fix $0<\epsilon\le1$. Suppose the actual monomial list for each original
    nonprivate gate approximates that gate in operator norm within
    $\delta_a=\epsilon/(16\max(1,M))$. Its monomials are allowed, each has at
    most $r$ pair effects, and its absolute coefficient sum is at most $B\ge1$.
    Rescale each such list by $(1+\delta_a)^{-1}$ to obtain a contraction
    circuit $W$. Set
    \begin{align}
        \delta_e & =\frac{\epsilon}{8\max(1,M)},                                 &
        m        & =\left\lceil\left(\frac{rB}{\delta_e}\right)^2\right\rceil+1.
    \end{align}
    Apply the constructed pair-effect replacement of length $m$ to $W$, and
    let $N$ be the number of sources of the resulting circuit. With
    \begin{align}
        Q & =B^{4b}D^4,                                                            &
        k & =\left\lceil\left(\frac{8\max(1,N)Q}{\epsilon}\right)^2\right\rceil+1,
    \end{align}
    there are normalized Schmidt data for these actual sources and a Gaussian
    sample of size $k$ whose sampled physical operator $\widehat\rho$ satisfies
    $\|\widehat\rho-\rho_G\|_1\le\epsilon$. Here $\rho_G$ is the physical
    density computed from the original gate operators, with all private
    registers traced out.
\end{theorem}
\begin{proof}
    \leanok
    \uses{thm:peps_sampled_original_density,thm:peps_approximate_physical_error,thm:peps_finite_whole_party_input}
    The input coordinates are obtained from finite-dimensionality of the
    original party memories. Rescaling preserves the original gate
    occurrences and participating sets, and does not increase the coefficient
    sums or pair-effect counts. The actual sampling theorem applied to $W$
    therefore gives $\|\widehat\rho-\rho_W\|_1\le3\epsilon/4$ with the
    stated replacement and sample counts. The local approximation estimate
    gives $\|\rho_W-\rho_G\|_1\le4M\delta_a\le\epsilon/4$. Thus
    \begin{align}
        \|\widehat\rho-\rho_G\|_1
         & \le\|\widehat\rho-\rho_W\|_1+\|\rho_W-\rho_G\|_1
        \le\epsilon.
    \end{align}
\end{proof}

\begin{theorem}[Explicit polynomial dimensions of the replacement links]
    \label{thm:peps_physical_first_link_polynomial}
    \lean{TNLean.PEPS.PairEffect.OriginalCircuit.card_physicalFirstNetworkAlphabet_le_power}
    \leanok
    Let an original circuit have at most $N$ nonprivate gate occurrences,
    each with at most $b$ participants. Suppose each gate has at most $K$
    monomials, at most $r$ pair effects in each monomial, and absolute
    coefficient sum at most $B$. Let $L\geq1$, and assume
    \begin{align}
        K\leq C_KL^u,\qquad N\leq C_NL^n,\qquad
        0\leq B\leq C_BL^s,\qquad 0\leq D\leq C_DL^d,
    \end{align}
    where $C_K,C_B\geq0$, $C_N\geq1$, and $u,n,s,d,p$ are nonnegative
    integers. Use the per-gate budget
    $\delta=L^{-p}/(8\max(1,N))$, and let $R$ be the resulting source
    replacement, with the physical output registers ordered first.
    For any nonnegative integer $\tau$, choose the actual sample count
    \begin{align}
        k=\left\lceil
        \bigl(8\max(1,\lvert\operatorname{Sources}(R)\rvert)
        B^{4\tau}D^4L^p\bigr)^2
        \right\rceil+1.
    \end{align}
    Define constants and exponents by
    \begin{align}
        C_{\mathrm{src}} & =\max\left(1,C_N\binom b2\right),                          \\
        C_{\mathrm{br}}  & =C_K\bigl(64C_N^2(rC_B)^2+2\bigr)^r,
                         & a_{\mathrm{br}}                      & =u+2r(n+s+p),       \\
        C_{\mathrm{sam}} & =64C_{\mathrm{src}}^2
        (C_B^{4\tau}C_D^4)^2+2,
                         & a_{\mathrm{sam}}                     & =2(n+4\tau s+4d+p).
    \end{align}
    For every retained set of gate occurrences, every root assignment,
    and every corresponding star or sample link $e$, its actual alphabet
    $\mathcal A_e$ satisfies
    \begin{align}
        \lvert\mathcal A_e\rvert
        \leq (C_{\mathrm{br}}^2+C_{\mathrm{sam}})
        L^{\max(2a_{\mathrm{br}},a_{\mathrm{sam}})}.
    \end{align}
    These constants depend only on the displayed uniform resource bounds.
    No private-memory dimension or source Schmidt rank occurs. This is a
    dimension estimate for the constructed alphabets; the error estimate
    and equality with a contraction of local tensors are separate assertions.
\end{theorem}
\begin{proof}
    \leanok
    \uses{thm:peps_physical_first_branch_polynomial,thm:peps_physical_first_sample_polynomial,thm:peps_source_alphabet_polynomial}
    Theorem~\ref{thm:peps_physical_first_branch_polynomial} gives at most
    $C_{\mathrm{br}}L^{a_{\mathrm{br}}}$ actual branches per gate.
    There are at most $N\binom b2$ actual source occurrences, so the
    estimate of Theorem~\ref{thm:peps_physical_first_sample_polynomial}
    gives $k<C_{\mathrm{sam}}L^{a_{\mathrm{sam}}}$.
    A star alphabet has cardinality equal to the square of its gate's
    branch count; a sample alphabet has cardinality $k$. Apply the
    bound of Theorem~\ref{thm:peps_source_alphabet_polynomial}
    to these two estimates. The use of
    $C_{\mathrm{src}}\geq1$ includes circuits without sources.
\end{proof}
